% Sujet en LaTeX
\documentclass{article}
\usepackage{hyperref}
\usepackage{geometry}
\usepackage{amsmath}
\geometry{a4paper, margin=1.5cm}

\title{Corrigé 1: Racines carrées et géométrie}
\author{}
\date{}

\begin{document}
\maketitle

\subsection*{Partie 1 : Rappel de la notion de racine carrée à travers les aires}

\begin{enumerate}
    \item \textbf{Construction de carrés :}
    \begin{enumerate}
        \item Les carrés ont les longueurs de côté suivantes :
        \begin{itemize}
            \item Pour une aire de $1~\text{cm}^2$, la longueur du côté est $\sqrt{1} = 1~\text{cm}$.
            \item Pour une aire de $4~\text{cm}^2$, la longueur du côté est $\sqrt{4} = 2~\text{cm}$.
            \item Pour une aire de $9~\text{cm}^2$, la longueur du côté est $\sqrt{9} = 3~\text{cm}$.
        \end{itemize}
    \end{enumerate}
    \item \textbf{Carrés d'aires non entières :}
    \begin{enumerate}
        \item Pour dessiner un carré d'aire $2~\text{cm}^2$, on construit un carré dont les côtés mesurent $\sqrt{2}~\text{cm}$.
        \item Construction :
        \begin{itemize}
            \item Tracez un segment de $1~\text{cm}$.
            \item À partir d'une extrémité, construisez un segment perpendiculaire de $1~\text{cm}$.
            \item Reliez les extrémités pour former un triangle rectangle.
            \item La longueur de l'hypoténuse sera $\sqrt{1^2 + 1^2} = \sqrt{2}~\text{cm}$.
            \item Utilisez cette longueur pour construire le carré d'aire $2~\text{cm}^2$.
        \end{itemize}
        \item La longueur du côté de ce carré est $\sqrt{2}~\text{cm}$.
    \end{enumerate}
\end{enumerate}

\subsection*{Partie 2 : Lois de multiplication des racines carrées}

\begin{enumerate}
    \item \textbf{Multiplication géométrique :}
    \begin{enumerate}
        \item L'aire du rectangle est :
        \[
        \sqrt{2} \times \sqrt{8} = \sqrt{16} = 4~\text{cm}^2.
        \]
    \end{enumerate}
    \item \textbf{Carré équivalent :}
    \begin{enumerate}
        \item La longueur du côté du carré équivalent est $\sqrt{4} = 2~\text{cm}$.
        \item Comparaison :
        \[
        \sqrt{2} \times \sqrt{8} = \sqrt{16} = 4 \quad \text{et} \quad \sqrt{2 \times 8} = \sqrt{16} = 4.
        \]
        \item Conclusion : $\sqrt{2} \times \sqrt{8} = \sqrt{16}$, ce qui illustre la loi de multiplication des racines carrées.
    \end{enumerate}
\end{enumerate}

\subsection*{Partie 3 : Exploration de l'inégalité $\sqrt{a + b} \leq \sqrt{a} + \sqrt{b}$}

\begin{enumerate}
    \item \textbf{Triangle rectangle :}
    \begin{enumerate}
        \item Avec $a = 9$ et $b = 16$ :
        \begin{itemize}
            \item $\sqrt{a} = \sqrt{9} = 3$.
            \item $\sqrt{b} = \sqrt{16} = 4$.
            \item $\sqrt{a + b} = \sqrt{25} = 5$.
        \end{itemize}
        \item Construction d'un triangle rectangle avec des côtés de $3~\text{cm}$ et $4~\text{cm}$.
        \item Longueur de l'hypoténuse :
        \[
        c = \sqrt{(\sqrt{a})^2 + (\sqrt{b})^2} = \sqrt{9 + 16} = \sqrt{25} = 5~\text{cm}.
        \]
    \end{enumerate}
    \item \textbf{Analyse de l'inégalité :}
    \begin{enumerate}
        \item Calcul de la somme :
        \[
        \sqrt{a} + \sqrt{b} = 3 + 4 = 7.
        \]
        \item Comparaison :
        \[
        \sqrt{a + b} = 5 \leq 7 = \sqrt{a} + \sqrt{b}.
        \]
        \item Conclusion : L'inégalité $\sqrt{a + b} \leq \sqrt{a} + \sqrt{b}$ est vérifiée pour $a = 9$ et $b = 16$.
    \end{enumerate}
\end{enumerate}

\end{document}
